\section{Variance structure}

\subsection{Setup}

Let $e_n$ denote the detrended gap and write the conditional decomposition
\[
e_n = \mu_n + u_n,\qquad \mathbb{E}(u_n \mid p_n \bmod M) = 0,
\]
where $\mu_n = \mathbb{E}(e_n \mid p_n \bmod M)$ is the phase-conditional mean and $u_n$ is the within-phase residual.

By the law of total variance,
\[
\operatorname{Var}(e) = \operatorname{Var}(\mu_n) + \mathbb{E}\!\left[\operatorname{Var}(u_n \mid r)\right].
\]
The variance split \texttt{var\_between} $+$ \texttt{var\_within} $=$ \texttt{var\_e} is reported in Table~\ref{tab:variance-split} (Section~3.3). This section interprets that split and records its relation to the covariance decomposition of Section~4.

\subsection{Monotone growth of the between-phase variance}

The between-phase variance \texttt{var\_between} grows monotonically with $M$ across the tested ladder, from $0.158$ at $M=6$ to $16.402$ at $M=9699690$. At the largest modulus this accounts for approximately $3.74\%$ of the total detrended variance $\operatorname{Var}(e) = 438.670412$. The within-phase variance declines correspondingly.

The growth is a purely arithmetic effect of the conditioning. As $M$ increases, more of the residue-class structure of the opening prime is resolved, and a larger fraction of the gap-size variance is attributed to phase. No claim is made that this fraction converges to a limit as $M\to\infty$; the seven moduli form a finite ladder.

\subsection{Relation to the covariance decomposition}

Section~4 decomposes the lag-one covariance as
\[
\operatorname{Cov}(e_n,\,e_{n+1}) = A + X + Y + D,
\]
where $A$ is the phase component, $X$ and $Y$ are cross terms, and $D$ is the residual term.

Both decompositions arise from the same phase conditioning. The between-phase variance \texttt{var\_between} and the phase covariance component $A$ capture the contribution of phase structure to two different second-order statistics:

\begin{itemize}
\item \texttt{var\_between} is a one-point statistic: it measures the dispersion of the phase-conditional means $\mu(r)$ around the global mean.
\item $A = \mathbb{E}[\mu(r_n)\mu(r_{n+1})]$ is a two-point statistic: it measures the lag-one autocovariance of the phase-conditional means.
\end{itemize}

The two quantities are not equivalent. A phase structure with large one-point dispersion need not produce a large lag-one autocovariance, and vice versa. The observed data show that both grow with $M$, but at different rates and with different signs in the small-$M$ regime.

For comparison, the non-$D$ covariance share $S(M)=(A+X+Y)/\operatorname{Cov}=1-D/\operatorname{Cov}$ increases from approximately $1.5\%$ at $M=6$ to $60.4\%$ at $M=9{,}699{,}690$ (Table~\ref{tab:abcd}). This should be read as an algebraic decomposition share, not as a causal fraction of dependence removed by conditioning. The $3.74\%$ between-phase variance share and the $60.4\%$ non-$D$ covariance share describe different one-point and two-point statistics and are not directly comparable effect sizes.

\subsection{What this section does not claim}

The variance split is an algebraic identity. It partitions the detrended variance into a between-phase and a within-phase term; it does not attribute a causal mechanism to either. The phase-conditional mean $\mu(r)$ is a computed quantity, not a modelled one.

No claim is made that the residual lag-one correlation $\rho_{\text{resid}}$ converges to zero, nor that the between-phase variance converges to a finite limit, nor that the phase decomposition explains all of the observed dependence. The empirical claims are those stated in Sections~4 and~5: the decomposition is exact, the residual remains negative at all seven tested moduli, and the pre-registered ranges are reproduced at $10^{10}$.
